\documentclass{article}
\usepackage{graphicx}
\usepackage{amsmath}
\usepackage{amsfonts}
\usepackage{bookmark}
\usepackage{booktabs}
\usepackage{exercises}
\graphicspath{ {} }
\usepackage{indentfirst}
\usepackage{amsthm}
\usepackage{yhmath}
\usepackage{geometry}
\usepackage{placeins}
\usepackage{csquotes}
\usepackage{pgfplots}
\newcommand*{\oldneg}{\mathord{\sim}}
\newcommand\showdiv[1]{\overline{\smash{)}#1}}
\newcommand\Mydiv[1]{%
\smash{\raise.3ex\hbox{$\big)$}}\mkern-8mu
        \overline{\enspace\strut#1}}
\setlength{\tabcolsep}{5pt}
\theoremstyle{definition}
\newtheorem{definition}{Definition}

\setlength{\parindent}{.5in}
\setlength{\parskip}{12pt}
\renewcommand{\arraystretch}{2}
\geometry{margin=1in}
\newcommand{\?}{\stackrel{?}{=}}
\pgfplotsset{width=7cm,compat=1.18}\usepgfplotslibrary{statistics}

\title{7.3 Answers}
\begin{document}
\maketitle

\section*{1}
The constant ratio of consecutive terms in a geometric sequence is called the common ratio.

\section*{2}
Suppose we are given a sequence $a_1,a_2,...,a_n$. One can determine if the sequence is geometric by checking if $a_n=a_1r^{n-1}$ for some common ratio $r$. This common ratio could be found 
by simply dividing $a_2$ by $a_1$.

\section*{12}
This is a geometric sequence, as it can be seen that each term in the sequence has a common ratio of $3$. In particular, the $n$th term of the sequence can be given by:
\[ a_n=0.2(3)^{n-1} \]

\section*{13}
This is not a geometric sequence, as there is no common ratio between each of the terms. In particular, while $\frac{a_2}{a_1}=\frac{10}{-5}=-2$, 
$\frac{a_3}{a_2}=\frac{20}{10}=2$.

\section*{14}
This is not a geometric sequence, as there is no common ratio between each of the terms. For example, while $\frac{a_2}{a_1}=\frac{1.5}{0.75}=2$,
$\frac{a_3}{a_2}=\frac{2.25}{1.5}=1.5$.

\section*{24}
The $n$th term of the sequence can be given by:
\[ a_n=7(-0.6)^{n-1} \]

\section*{25}
The $n$th term of the sequence can be given by:
\[ a_n=5(-2.8)^{n-1} \]

\section*{26}
The $n$th term of the sequence can be given by:
\[ a_n=120(1.5)^{n-1} \]

\section*{27}
B

\section*{37}
The error here is that $2$ was raised to the $n$th power, rather than the $n-1$th power. Instead, the answer is $a_n=3(2)^{n-1}$.

\section*{45}
We know that \[ 162=a_1r^3 \] and \[ 4374=a_1r^6 \] meaning \[ a_1=\frac{162}{r^3} \] and thus \[ 4374=a_1r^6 \] \[ 4374=\frac{162}{r^3}r^6 \]
\[ 4374=162r^3 \] \[ 27=r^3 \] \[ r=3 \] and $a_1=\frac{162}{27}=6$ meaning the rule for the $n$th term of the sequence is $a_n=6(3)^{n-1}$.

\section*{46}
We know that \[ \frac{7}{4}=a_1r^2 \] and \[ \frac{7}{16}=a_1r^4 \] meaning \[ a_1=\frac{7}{4r^2} \] and thus \[ \frac{7}{16}=a_1r^4 \]
\[ \frac{7}{16}=\frac{7}{4r^2}r^4 \] \[ \frac{28}{16}=7r^2 \] \[ r^2=\frac{4}{16} \] \[ r=\pm\frac{1}{2} \] Thus we know that $a_1=\frac{7}{4\frac{1}{4}}=7$. 
Thus the rule for the sequence is either $a_n=7(\frac{1}{2})^{n-1}$ or $a_n=7(-\frac{1}{2})^{n-1}$.

\section*{47}
(My answer is not getting the correct value for $a_7$, but I cannot find any errors in the solution)\newline
We know that \[ 6=a_1r^3 \] and \[ \frac{243}{8}=a_1r^6 \] meaning \[ a_1=\frac{6}{r^3} \] and thus \[ \frac{243}{8}=a_1r^6 \]
\[ \frac{243}{8}=\frac{6}{r^3}r^6 \] \[ \frac{243}{48}=r^3 \] \[ r=\frac{3\sqrt[3]{9}}{2\sqrt[3]{6}}=\frac{3\sqrt[3]{9\times 36}}{2} \]
\[ =\frac{9\sqrt[3]{12}}{2} \] and thus \[ a_1=\frac{6}{(\frac{9\sqrt[3]{12}}{2})^3}=\frac{48}{8748}=\frac{12}{2187} \] meaning
the rule for the sequence is $a_n=\frac{12}{2187}(\frac{9\sqrt[3]{12}}{2})^{n-1}$.

\section*{51}
We know that $a_1=4$ and $r=\frac{1}{4}$, meaning the sum of the first 6 terms of the sequence is $4\frac{1-(\frac{1}{4})^6}{1-\frac{1}{4}}=\frac{1365}{256}$.

\section*{52}
We know that $a_1=8$ and $r=\frac{3}{2}$, meaning the sum of the first 12 terms of the sequence is $8\frac{1-(\frac{3}{2})^{12}}{1-\frac{3}{2}}=\frac{527345}{256}$.

\section*{53}
We know that $a_1=1$ and $r=-4$, meaning the sum of the first 11 terms of the sequence is $\frac{1-(-4)^{11}}{1+4}=838861$.

\section*{54}
We know that the $a_1=2$ and $r=3$, meaning the sum of the first $9$ terms of the sequence is $2\frac{1-3^9}{1-3}=19682$. C

\section*{56}
(a): We are given a geometric series with $a_1=1$ and $r=x$. Thus, the sum of the first 5 terms (which is the given polynomial) is 
$\frac{1-x^5}{1-x}$.


(b): We are given a geometric series with $a_1=3x$ and $r=2x^2$. Thus, the sum of the first 4 terms is
$\frac{1-(2x^2)^4}{1-2x^2}=\frac{1-16x^8}{1-2x^2}$.

\section*{59}
(a): We are given, after $0$ passes, an initial setup of 1024 items, and after each succeeding pass, that amount decreases by half. Thus, 
the rule is $a_n=1024(\frac{1}{2})^n$ where $a_0=1024$.


(b): We are trying to find what value of $n$ makes $a_n=1$. Thus \[ 1=1024(\frac{1}{2})^n \] \[ \frac{1}{1024}=(\frac{1}{2})^n \]
\[ \log_{\frac{1}{2}}\frac{1}{1024}=n \] \[ n=\frac{\log\frac{1}{1024}}{\log\frac{1}{2}}=10 \]

\section*{60}
(a): We are given that $a_1=1$, and $r=8$. Thus the rule is $a_n=8^{n-1}$. The total number of squares removed through stage 8 is the sum of the series of the first 8 terms 
in the sequence, which is $\frac{1-8^8}{1-8}=2396745$.


(b): The remaining area of the square at $a_1$ is $\frac{8}{9}$. At $a_2$, it is $\frac{64}{81}$. At $a_3$, it is $\frac{512}{729}$. Thus, the common ratio 
$r$ in this pattern is $\frac{8}{9}$, meaning the rule is $a_n=\frac{8}{9}(\frac{8}{9})^{n-1}$.

\section*{61}
\begin{samepage}
We are given that the initial amount is $\$2000$, and the annual interest rate is $5\%$, and $\$2000$ more is deposited at the start of each year. Let us consider an example of 
the total amount in the account immediately after despositing the \$2000 for the fifth time. This will be:
\[ (((2000(1.05)+2000)1.05+2000)1.05+2000)1.05+2000 \]
\[ =(2000+\frac{2000}{1.05}+\frac{2000}{(1.05)^2}+\frac{2000}{(1.05)^3})(1.05)^4+2000 \]
Note that \[ 2000+\frac{2000}{1.05}+\frac{2000}{(1.05)^2}+\frac{2000}{(1.05)^3} \] is the sum of a geometric series, where $a_1=2000$, and $r=\frac{1}{1.05}$, meaning 
$a_n=2000(\frac{1}{1.05})^{n-1}$. Now, let $s_n$ be the sum of that geometric series for its first $n$ terms. It is
\[ 2000\left(\dfrac{1-(\dfrac{1}{1.05})^n}{1-\dfrac{1}{1.05}}\right) \] 
Extending this pattern more generally, we can see that the amount in the account immediately after the $n$th deposit is equal to:
\[ t_n=s_{n-1}(1.05)^{n-1}+2000=2000\left(\dfrac{1-(\dfrac{1}{1.05})^{n-1}}{1-\dfrac{1}{1.05}}\right)(1.05)^{n-1}+2000 \]
where $t_1=2000$. Thus, immediately after making the 30th deposit, the account would have:
\[ 2000\left(\dfrac{1-(\dfrac{1}{1.05})^{30-1}}{1-\dfrac{1}{1.05}}\right)(1.05)^{30-1}+2000\approx \$132877.70 \]
\end{samepage}

\end{document}
